% !TeX root = ../../Book3.tex
% 纲要标题：## 第17章　投射维数与 Auslander–Buchsbaum 公式
% 纲要位置：计划纲要.md，第 2488 行。

\chapter{投射维数与 Auslander–Buchsbaum 公式}
\label{b3:ch:17}
\BookChapterNumber{b3:ch:17}

\begin{chapterabstract}
交换 Noether 局部环上的有限生成模，可以按照极小生成元逐步构造自由消解。在极小自由消解中，所有正次微分矩阵的系数都落在极大理想中，因而模掉极大理想后微分全为零；消解各项的秩于是成为可计算的 Betti 数。它们检测投射维数以及消解能否终止。对非零有限生成模，有限投射维数与深度由 Auslander–Buchsbaum 公式联系起来；周期消解则显示有限性条件为何必要。
\end{chapterabstract}

\begin{learninggoals}
\begin{itemize}
\item 构造局部极小自由消解，区分极小消解长度与任意消解长度。
\item 用剩余域的 Tor 或 Ext 计算 Betti 数，并检验有限生成模的投射维数。
\item 沿短正合列作维数平移，说明除以共同非零因子后的消解变化。
\item 证明并使用 Auslander–Buchsbaum 公式，识别无限投射维数造成的失效情形。
\end{itemize}
\end{learninggoals}

设 \(k\) 为域，令 \(R=k[x]/(x^2)\)，其剩余域自然识别为 \(k\)。乘以 \(x\) 的核恰是理想 \((x)\)，也是它的像，于是得到不断重复的正合列
\[
\cdots\xrightarrow{\,x\,}R\xrightarrow{\,x\,}R\xrightarrow{\,x\,}R\longrightarrow k\longrightarrow0.
\]
每一项都是自由模，微分却永远不能删去；模掉极大理想后，所有微分都变为零。这个例子说明自由消解可以无穷延伸，而各层所需的生成元数仍有稳定含义。

\ProseParagraphBreak

换成一般 Noether 局部环，若某个自由消解中的微分含单位系数，就可以拆掉一对相邻的自由生成元。这对生成元组成一个可缩的二项复形，去掉它并不改变同调。当消解长度和各项的秩都有限时，只需有限次这样的操作，就能得到极小自由消解。对可能无限长的消解，则可以直接按各层关系的极小生成组逐步构造。所得极小消解的秩是模自身的数据；它何时终止，便是投射维数问题。深度则从正则序列的长度衡量局部环和模的性质，有限投射维数时二者由 Auslander–Buchsbaum 公式精确联系。

\ProseParagraphBreak

本章使用\cref{b3:cor:local-generators}、\cref{b3:thm:koszul-acyclic}与\cref{b3:thm:depth-ext}，并沿用前章已给出的 Ext、Tor 基础。自由、投射、张量积与 Hom 来自第二卷第3、6、7章；一般消解比较及维数平移的系统论述安排在第四卷第6—9章，其中第9.5节专门处理维数平移。当前所需的具体形式已在本卷正文证明。

\ProseParagraphBreak
极小消解及末端 Ext 的非消失可参阅松村英之的《Commutative Ring Theory》\cite[\S19, Minimal free resolutions, Lemma 1, pp. 153--154]{Matsumura1989CommutativeRingTheory}；本章中心公式见同书\cite[Theorem 19.1, p. 155]{Matsumura1989CommutativeRingTheory}。

% !TeX root = ../../Book3.tex
% 纲要标题：### 17.1 局部极小自由消解
% 纲要位置：计划纲要.md，第 2490 行。

\section{局部极小自由消解}
\label{b3:sec:1701}

自由消解中的一个单位系数把相邻两项的两个基向量配成一对，可以通过换基把这一对单独拆出。极小自由消解排除这样的配对，要求所有正次微分的矩阵元素都落在局部环的极大理想中。约化到剩余域后，微分同时变成零，各阶自由秩便能从模的同调中读取。

% 纲要内容（仅作写作计划，尚非教材正文）：
%
% * 极小生成元与极小微分
% * 单位系数消去及有限长度、各项有限秩消解的可缩直和分解
% * 剩余域张量后的简化
% * Betti 数及终止性
%

\subsection{极小生成元与极小微分}
\label{b3:subsec:1701:01}

给有限生成模写一个自由满射并不困难；困难在于分清哪些生成元确有必要，以及关系的关系会不会一直出现。局部环的剩余域把这两个问题变成向量空间问题。本节的自由消解各项均为有限秩自由模，但不预设消解长度有限，仍采用降次链复形记号。

\begin{definition}\label{b3:def:minimal-free-resolution}
设 \((R,\mathfrak m)\) 为交换 Noether 局部环，\(k=R/\mathfrak m\) 为其剩余域，\(M\) 为有限生成的 \(R\) 模。若增广正合复形
\[
\cdots\longrightarrow F_2\xrightarrow{\mathrm{d}_2}F_1
\xrightarrow{\mathrm{d}_1}F_0\xrightarrow{\epsilon}M\longrightarrow0
\]
中每个 \(F_i\) 为有限秩自由模，且 \(\mathrm{d}_i(F_i)\subseteq\mathfrak mF_{i-1}\) 对 \(i\ge1\) 成立，则称它为 \(M\) 的\term[jixiao-ziyou-xiaojie]{极小自由消解}[minimal free resolution]。
\end{definition}

正合性与 \(\mathrm{d}_1(F_1)\subseteq\mathfrak mF_0\) 说明增广映射诱导同构
\[
F_0/\mathfrak mF_0\xrightarrow{\sim}M/\mathfrak mM.
\]
因此极小性也包括零次项没有多余生成元。矩阵表示下，条件就是所有正次微分的矩阵元素都属于 \(\mathfrak m\)。

\ProseParagraphBreak
设 \(R\) 为交换环。若一个 \(R\) 模链复形的恒等链映射与零映射链同伦，则称它为\term[kesuo-fuxing]{可缩复形}[contractible complex]。相邻两次的恒等映射 \(R\xrightarrow{1}R\) 是最简单的例子；下面的命题说明它怎样从含单位系数的微分中出现。

\begin{proposition}[单位系数的消去]\label{b3:prop:unit-cancellation}
设 \((R,\mathfrak m)\) 为交换 Noether 局部环，\(k=R/\mathfrak m\)，\(M\) 为有限生成的 \(R\) 模，\(F_\bullet\to M\) 为各项均为有限秩自由模的自由消解。若某个整数 \(j\ge1\) 的微分 \(\mathrm{d}_j:F_j\to F_{j-1}\) 的矩阵含单位系数，则通过换基有增广复形的直和分解
\[
F_\bullet\cong G_\bullet\oplus\mathbb D(j),\qquad
\mathbb D(j):\quad0\longrightarrow R\xrightarrow{1}R\longrightarrow0,
\]
其中 \(\mathbb D(j)\) 的两个非零项位于次数 \(j,j-1\)，增广为零，它为可缩复形；\(G_\bullet\to M\) 仍为各项均为有限秩自由模的自由消解。由此，\(F_\bullet\) 极小等价于 \(F_\bullet\otimes_Rk\) 的全部正次微分为零，也等价于不能拆出任何 \(\mathbb D(j)\) 这样的非零直和项。
\end{proposition}
\begin{proof}
把单位系数交换到矩阵左上角，乘以其逆将它化为 \(1\)，再用可逆行、列变换清除同行、同列的其他系数，得到 \(\operatorname{diag}(1,D)\)。因此可选分解
\[
\begin{aligned}
F_j&=Ru\oplus F'_j,& F_{j-1}&=Rv\oplus F'_{j-1},\\
\mathrm{d}_j(u)&=v,& \mathrm{d}_j(F'_j)&\subseteq F'_{j-1}.
\end{aligned}
\]
由 \(\mathrm{d}_j\mathrm{d}_{j+1}=0\)，\(\mathrm{d}_{j+1}\) 的像在 \(u\) 方向的坐标必须为零。若 \(j>1\)，\(\mathrm{d}_{j-1}\mathrm{d}_j=0\) 给出 \(\mathrm{d}_{j-1}(v)=0\)；若 \(j=1\)，则由 \(\epsilon\mathrm{d}_1=0\) 得到 \(\epsilon(v)=0\)。所以 \(Ru\xrightarrow{1}Rv\) 与其余各项分离，给出所述直和分解。取 \(h_{j-1}(v)=u\)，其余同伦分量为零，便有 \(\mathrm{d}h+h\mathrm{d}=\operatorname{id}_{\mathbb D(j)}\)。此两项复形正合，因此余下的增广复形 \(G_\bullet\to M\) 也正合。

张量剩余域后微分为零，恰好表示原微分的全部矩阵系数属于 \(\mathfrak m\)。局部环中 \(\mathfrak m\) 之外的元素都是单位，所以不极小时可以作上述消去；极小时则在任何基下都不能出现单位系数，也不能有这样的两项直和因子。
\end{proof}

\begin{corollary}\label{b3:cor:finite-resolution-splitting}
设 \((R,\mathfrak m)\) 为交换 Noether 局部环，\(M\) 为有限生成的 \(R\) 模。给定 \(M\) 的有限长度、各项均为有限秩自由模的自由消解 \(F_\bullet\to M\)，可经有限次单位系数消去得到
\[
F_\bullet\cong F_\bullet^{\min}\oplus
\bigoplus_{\nu=1}^s\mathbb D(j_\nu),
\]
其中 \(F_\bullet^{\min}\to M\) 是极小自由消解，\(s\) 为非负整数，每个 \(\mathbb D(j_\nu)\) 是在次数 \(j_\nu,j_\nu-1\) 放置 \(R\xrightarrow{1}R\) 的两项复形。
\end{corollary}
\begin{proof}
原消解只有有限多个非零项，各项秩又有限，故总自由秩为有限整数。每次应用\cref{b3:prop:unit-cancellation}使总自由秩减少二，所以只需有限次消去。停止后所有正次微分均没有单位系数，即其系数全在 \(\mathfrak m\) 中，余下的消解正是极小自由消解。
\end{proof}

\begin{proposition}\label{b3:prop:minimal-resolution-existence}
设 \((R,\mathfrak m)\) 为交换 Noether 局部环，\(k=R/\mathfrak m\)，\(M\) 为有限生成的 \(R\) 模。则 \(M\) 有极小自由消解。它可通过依次选取满射 \(F_i\to K_i\) 构造，其中 \(K_0=M\)、\(K_{i+1}=\ker(F_i\to K_i)\)，且每个 \(F_i/\mathfrak mF_i\to K_i/\mathfrak mK_i\) 都是 \(k\) 线性同构。
\end{proposition}
\begin{proof}
取 \(M/\mathfrak mM\) 的一组 \(k\) 基并提升到 \(M\)。由\cref{b3:cor:local-generators}，提升构成极小生成组，给出 \(F_0\to M\)。因 \(R\) Noether 且 \(F_0\) 有限生成，核 \(K_1\) 也有限生成。对 \(K_1\) 重复这一过程，得到 \(F_1\to K_1\)，再合成到 \(F_0\)。

一般地令 \(K_0=M\)，依次取按 \(K_i/\mathfrak mK_i\) 的一组基提升所得的极小满射 \(\pi_i:F_i\to K_i\)，核为 \(K_{i+1}\)。所选基给出同构
\[
\overline{\pi_i}:F_i/\mathfrak mF_i\xrightarrow{\sim}K_i/\mathfrak mK_i.
\]
若 \(z\in K_{i+1}\)，则 \(\pi_i(z)=0\)，故 \(\overline{\pi_i}(z+\mathfrak mF_i)=0\)；由上式的单射性得 \(z\in\mathfrak mF_i\)。因此 \(K_{i+1}\subseteq\mathfrak mF_i\)，接起这些短正合列后，每个正次微分的像都属于下一项的 \(\mathfrak m\) 倍，得到所需极小消解。若某个核为零，以后全取零项。
\end{proof}

例如，设 \(k\) 为域，在 \(R=k[x,y]_{(x,y)}\) 上，模 \(R/(x,y)\) 有极小消解
\[
0\longrightarrow R\xrightarrow{\binom{-y}{x}}R^2
\xrightarrow{(x\quad y)}R\longrightarrow k\longrightarrow0.
\]
这也是已证的二元正则序列 Koszul 消解。在任何相邻的非负次数额外直加 \(R\xrightarrow{1}R\)，仍得到一个自由消解。这种加入可缩直和项的操作与\cref{b3:prop:unit-cancellation}中的消去互为逆过程。

\subsection{剩余域张量后的简化}
\label{b3:subsec:1701:02}

设 \(R\) 为交换环，\(M,N\) 为 \(R\) 模。本章所需的 Tor、Ext 来自自由消解。具体地，取 \(M\) 的自由消解 \(F_\bullet\)，约定
\[
\operatorname{Tor}_i^R(M,N)=H_i(F_\bullet\otimes_RN),\qquad
\operatorname{Ext}_R^i(M,N)=H^i\operatorname{Hom}_R(F_\bullet,N).
\]
由\cref{b3:lem:free-resolution-comparison}，消解之间的比较映射唯一到链同伦，故以上同调不依赖消解的选择。\cref{b3:lem:ext-basic-properties,b3:lem:tor-basic-properties}给出长正合列及 Tor 的平衡计算。对于极小消解，微分的像落在极大理想乘以相应自由模之中；取剩余域系数后，微分全部变为零，同调因而直接由消解各项得到。

\begin{proposition}\label{b3:prop:minimal-resolution-betti}
设 \((R,\mathfrak m)\) 为交换 Noether 局部环，\(k=R/\mathfrak m\)，\(M\) 为有限生成的 \(R\) 模，\(F_\bullet\to M\) 为其极小自由消解。则对每个整数 \(i\ge0\)，
\[
\operatorname{Tor}_i^R(M,k)\cong F_i\otimes_Rk,
\qquad
\operatorname{Ext}_R^i(M,k)\cong\operatorname{Hom}_R(F_i,k).
\]
任意两个极小自由消解同构为增广复形。
\end{proposition}
\begin{proof}
微分矩阵的元素在 \(\mathfrak m\) 中，张量 \(k\) 或映入 \(k\) 后全部成为零，所以同调就是相应项本身。
设 \(F_\bullet,F'_\bullet\) 为两个极小消解。由\cref{b3:lem:free-resolution-comparison}，取提升 \(M\) 的恒等映射的双向比较映射 \(f:F_\bullet\to F'_\bullet\)、\(g:F'_\bullet\to F_\bullet\)。两方向的复合各自与恒等链同伦；模掉 \(\mathfrak m\) 后，微分为零，同伦公式的右边也为零。因此
\[
\bar g_i\bar f_i=\operatorname{id}_{F_i/\mathfrak mF_i},
\qquad
\bar f_i\bar g_i=\operatorname{id}_{F'_i/\mathfrak mF'_i}.
\]
所以 \(\bar f_i:F_i/\mathfrak mF_i\to F'_i/\mathfrak mF'_i\) 是 \(k\) 线性同构，两边的自由秩相等。若共同秩为零，\(f_i\) 已是零模之间的同构；否则选基后 \(f_i\) 是方阵，且 \(\det(\bar f_i)\ne0\)，于是 \(\det(f_i)\notin\mathfrak m\)。局部环中极大理想之外的元素为单位，伴随矩阵公式因而给出 \(f_i\) 的逆。各次 \(f_i\) 可逆，且与微分和增广相容，构成增广复形同构。
\end{proof}

\subsection{Betti 数与消解终止性}
\label{b3:subsec:1701:03}

\begin{definition}\label{b3:def:local-betti-numbers}
设 \((R,\mathfrak m)\) 为交换 Noether 局部环，\(k=R/\mathfrak m\)，\(M\) 为有限生成的 \(R\) 模，\(i\ge0\) 为整数。整数
\(\beta_i^R(M)\coloneqq\dim_k\operatorname{Tor}_i^R(M,k)\)
称为 \(M\) 的第 \(i\) 个\term[betti-shu]{Betti 数}[Betti number]。
\end{definition}

\(\beta_0\) 是最少生成元数，\(\beta_1\) 是极小展示中的最少关系数。更高的数统计关系之间的关系。它们依赖底环：同一个剩余域 \(k\) 作为 \(k\) 模只有 \(\beta_0=1\)，作为 \(k[x,y]_{(x,y)}\) 模则有 \((\beta_0,\beta_1,\beta_2)=(1,2,1)\)。

\begin{corollary}\label{b3:cor:betti-gap-termination}
设 \((R,\mathfrak m)\) 为交换 Noether 局部环，\(k=R/\mathfrak m\)，\(M\) 为有限生成的 \(R\) 模，\(i\ge0\) 为整数。若 \(\beta_i^R(M)=0\)，则 \(M\) 的极小自由消解从第 \(i\) 项起全部为零。\(i=0\) 时，这一条件等价于 \(M=0\)。
\end{corollary}
\begin{proof}
\cref{b3:prop:minimal-resolution-betti}给出 \(F_i=0\)。若 \(i>0\)，极小满射 \(F_i\to K_i\) 使 \(K_i=0\)，后续所有核和极小自由项也为零。\(i=0\) 时满射 \(F_0\to M\) 直接给出 \(M=0\)。
\end{proof}

作为对照，设 \(k\) 为域，令 \(R=k[\epsilon]/(\epsilon^2)\)。其剩余域具有消解
\[
\cdots\xrightarrow{\epsilon}R\xrightarrow{\epsilon}R
\xrightarrow{\epsilon}R\longrightarrow k\longrightarrow0.
\]
乘 \(\epsilon\) 的核与像均为 \((\epsilon)\)，所以它正合且极小；每个 Betti 数均为一。有限生成并不保证消解终止。

% !TeX root = ../../Book3.tex
% 纲要标题：### 17.2 投射维数
% 纲要位置：工作记录/计划纲要.md，第 2526 行。
% 纲要细目（写作范围）：
% * 有限生成模的投射维数
% * Ext、Tor 与投射维数判据
% * 维数平移及一般投射模的 Ext 判据在此证明；第四卷第9.5节与第11章作系统整理
% * 剩余域上界、正则元换环及极小消解的同调对照表
% * 测试模选择的 Tor、Ext 反例及完备化保持 Betti 数、投射维数

\section{投射维数}
\label{b3:sec:1702}

一个模有无限自由消解，并不表示每种消解都必须无限长。投射维数问最短长度，而 Tor 消失与极小消解的末项提供可以实际检测这一长度的条件。

% 纲要内容（仅作写作计划，尚非教材正文）：
%
% * 有限生成模的投射维数
% * Ext、Tor 与投射维数判据
% * 维数平移及一般投射模的 Ext 判据在此证明；第四卷第9.5节与第11章作系统整理
% * 剩余域上界、正则元换环及极小消解的同调对照表
%

\subsection{有限生成模的投射维数}
\label{b3:subsec:1702:01}

\begin{definition}\label{b3:def:projective-dimension}
设 \(R\) 为交换环，\(M\ne0\) 为 \(R\) 模。若存在正合列
\[
0\longrightarrow P_n\longrightarrow\cdots\longrightarrow P_0
\longrightarrow M\longrightarrow0
\]
且各 \(P_i\) 投射，则使这种列存在的最小非负整数 \(n\) 称为 \(M\) 的\term[toushe-weishu]{投射维数}[projective dimension]，记 \(\operatorname{pd}_RM\)。若不存在，记为 \(\infty\)；约定 \(\operatorname{pd}_R0=-\infty\)。
\end{definition}

定义中的 \(P_i\) 只要求投射，并不预设自由或有限生成。对 Noether 局部环上的有限生成模，极小自由消解的各项虽都有限自由，却已经实现这类投射消解的最短长度。投射维数衡量消解必须延伸多少步；任意消解却可以插入 \(R\xrightarrow{1}R\)，人为增加长度。

\begin{proposition}\label{b3:prop:pd-minimal-length}
设 \((R,\mathfrak m,k)\) 为交换 Noether 局部环，\(M\ne0\) 为有限生成的 \(R\) 模。若 \(F_\bullet\rightarrow M\) 是极小自由消解，则
\[
\operatorname{pd}_RM=\sup\{\,i:F_i\ne0\,\}
=\sup\{\,i:\operatorname{Tor}_i^R(M,k)\ne0\,\}.
\]
\end{proposition}
\begin{proof}
投射模的提升性质使自由消解的比较映射论证同样适用于投射消解。因此有限投射消解在高于其长度的次数计算出零 Tor，故极小消解的这些项由\cref{b3:prop:minimal-resolution-betti}为零。反向，若极小消解在次数 \(r\) 的项非零，而 \(F_{r+1}=0\)，则全部后续项也为零，它本身就是长度 \(r\) 的投射消解；前一方向又排除了任何更短的投射消解，故两个上确界都为 \(r\)。若每个 \(F_i\) 都非零，则 Tor 在任意高次数均非零，不可能另有有限投射消解，两个上确界及投射维数都为 \(\infty\)。
\end{proof}

\subsection{Ext、Tor 与投射维数判据}
\label{b3:subsec:1702:02}

极小消解的微分在剩余域 \(k\) 上全为零，所以 \(\operatorname{Tor}_i^R(M,k)\) 和 \(\operatorname{Ext}_R^i(M,k)\) 直接记录 \(F_i\) 的秩。任一项同调为零就迫使该自由项为零，\cref{b3:cor:betti-gap-termination}再使所有后续项为零。因此在这里，仅测试剩余域便能判断消解是否终止，无须逐个测试所有模。

\begin{theorem}\label{b3:thm:local-pd-vanishing}
设 \((R,\mathfrak m,k)\) 为交换 Noether 局部环，\(M\) 为有限生成的 \(R\) 模，\(n\ge0\) 为整数。下列条件等价：
\begin{equivalences}
\item \(\operatorname{pd}_RM\le n\)；
\item \(\operatorname{Tor}_{n+1}^R(M,k)=0\)；
\item \(\operatorname{Ext}_R^{n+1}(M,k)=0\)；
\item 对每个 \(R\) 模 \(N\)，有 \(\operatorname{Ext}_R^{n+1}(M,N)=0\)。
\end{equivalences}
若 \(M\ne0\) 且 \(r=\operatorname{pd}_RM<\infty\)，则每个非零有限生成的 \(R\) 模 \(N\) 都满足 \(\operatorname{Ext}_R^r(M,N)\ne0\)。
\end{theorem}
\begin{proof}
\(M=0\) 时所有条件直接成立。设 \(M\ne0\)，前面三项分别等价于极小消解的 \(F_{n+1}=0\)，由\cref{b3:cor:betti-gap-termination}得到同一终止条件。第一项使 Hom 复形在次数 \(n+1\) 后为零，推出第四项；第四项取 \(N=k\) 推出第三项。

若最后非零项为 \(F_r\)，约定 \(r=0\) 时下面的像为零，便可统一写成
\[
\operatorname{Ext}_R^r(M,N)
=\operatorname{Hom}_R(F_r,N)/\operatorname{im}(\mathrm{d}_r^*).
\]
极小性给 \(\operatorname{im}(\mathrm{d}_r^*)\subseteq\mathfrak m\operatorname{Hom}_R(F_r,N)\)。记 \(T=\operatorname{Hom}_R(F_r,N)\)，它是非零有限个 \(N\) 的直和，所以非零且有限生成。对局部环 \(R\) 上的 \(T\) 用 Nakayama 引理，得到 \(T/\mathfrak mT\ne0\)。所列 Ext 商仍满射到 \(T/\mathfrak mT\)，故非零；这里必须保留 \(N\) 非零且有限生成的条件。\(r=0\) 时像为零，同样成立。
\end{proof}

任意其他测试模未必能看见这些自由项；\cref{b3:prop:pd-test-module-counterexamples}将给出 Tor 的测试模换成一个非零商、Ext 的测试模换成环本身后判据失效的计算。

\ProseParagraphBreak
对非零有限生成局部模，深度取 \(\operatorname{Ext}_R^i(k,M)\) 的首次非零次数，投射维数则取 \(\operatorname{Ext}_R^i(M,k)\) 的非零次数上确界；后者可能为无穷。两个变量和两个端点都不同。在 \(R=k[X,Y]_{(X,Y)}\) 上，即使两边都取同一个模 \(M=k\)，所得数值仍不同。用 \(X,Y\) 的 Koszul 消解映入 \(k\)，微分全为零，零、一、二次 Ext 的 \(k\) 维数分别为 \(1,2,1\)，其余为零。因此
\[
\begin{array}{c|c|c}
\text{检测目标}&\text{Ext 计算}&\text{所得数值}\\\hline
\operatorname{depth}_Rk&\operatorname{Ext}_R^0(k,k)=k,
\ \text{首次非零次数为 }0&0\\
\operatorname{pd}_Rk&\operatorname{Ext}_R^2(k,k)=k,
\ \operatorname{Ext}_R^{i>2}(k,k)=0&2
\end{array}
\]
同一组 Ext 的首次非零和最后非零分别给出深度与投射维数；把两个端点混为一谈，也会在这个例子中产生错误。

\subsection{维数平移}
\label{b3:subsec:1702:03}

这里的“维数平移”指同调次数的平移：从一个模换到它的关系模，Ext 和 Tor 的次数相应改变一，与 Krull 维数 \(\dim R\) 无关。

\begin{proposition}\label{b3:prop:dimension-shifting-local}
设 \(R\) 为交换环，\(0\rightarrow K\rightarrow P\rightarrow M\rightarrow0\) 正合且 \(P\) 投射。对任意 \(R\) 模 \(N\) 和 \(i\ge1\)，有自然同构
\[
\operatorname{Ext}_R^i(K,N)\cong\operatorname{Ext}_R^{i+1}(M,N),
\qquad
\operatorname{Tor}_i^R(K,N)\cong\operatorname{Tor}_{i+1}^R(M,N).
\]
\end{proposition}
\begin{proof}
投射模是自由模的直和项，故它的所有正次 Ext 和 Tor 为零。在\cref{b3:lem:ext-basic-properties}和\cref{b3:lem:tor-basic-properties}的长正合列中，连接映射两边相邻的 \(P\) 项都为零，因而连接映射为同构。
\end{proof}

\begin{proposition}[投射模的 Ext 判据]\label{b3:prop:projective-ext-criterion}
设 \(R\) 为交换环，\(P\) 为 \(R\) 模。以下条件等价：
\begin{equivalences}
\item \(P\) 投射；
\item 对每个 \(R\) 模 \(N\)，有 \(\operatorname{Ext}_R^1(P,N)=0\)。
\end{equivalences}
\end{proposition}
\begin{proof}
若 \(P\) 投射，则 \(P\) 是某个自由模 \(F\) 的直和项。Ext 在第一变量保持有限直和，而自由模的正次 Ext 为零，因此 \(\operatorname{Ext}_R^1(P,N)=0\) 对每个 \(N\) 成立。

反过来，取自由模满射 \(q:F\rightarrow P\)，令 \(K=\ker q\)。要把 \(P\) 实现为自由模的直和项，只需找到截面 \(s:P\to F\)，使 \(qs=\operatorname{id}_P\)；这就是把恒等映射沿 \(q\) 提升。对短正合列 \(0\rightarrow K\rightarrow F\xrightarrow{q}P\rightarrow0\) 应用第二变量的 Ext 长正合列，得到
\[
\operatorname{Hom}_R(P,F)\xrightarrow{q\circ-}\operatorname{Hom}_R(P,P)
\longrightarrow\operatorname{Ext}_R^1(P,K)=0.
\]
因此所需的 \(s\) 存在。每个 \(u\in F\) 都可写成 \(sq(u)+(u-sq(u))\)，第二项属于 \(K\)，而 \(s(P)\cap K=0\)，所以 \(F=s(P)\oplus K\)。原短正合列分裂，\(P\) 是自由模 \(F\) 的直和项，故投射。
\end{proof}

设 \((R,\mathfrak m)\) 为交换 Noether 局部环，\(M\) 为有限生成的 \(R\) 模。令 \(\Omega^0M=M\)，并令 \(\Omega^{j+1}M\) 为极小满射 \(F_j\rightarrow\Omega^jM\) 的核；这些是\cref{b3:def:syzygy-local}中合冲模在极小消解上的选择。沿消解换到下一层合冲模，每次便将相应 Ext 的次数降低一。迭代 \(n\) 次后，对 \(n\ge0\) 与任意 \(R\) 模 \(N\)，有
\[
\operatorname{Ext}_R^1(\Omega^nM,N)\cong\operatorname{Ext}_R^{n+1}(M,N).
\]

特别地，若 \(M\) 的投射维数为有限正整数 \(r\)，则 \(\Omega M\ne0\) 且
\(\operatorname{pd}_R\Omega M=r-1\)。极小消解去掉零次项就是 \(\Omega M\) 的极小消解，所以这个等式还可以不经一般判据而直接看出。

\subsection{剩余域上界与正则元换环}
\label{b3:subsec:1702:04}

\begin{corollary}\label{b3:cor:residue-pd-bound}
设 \((R,\mathfrak m,k)\) 为交换 Noether 局部环。每个有限生成的 \(R\) 模 \(M\) 都满足
\(\operatorname{pd}_RM\le\operatorname{pd}_Rk\)。有限生成模的直和因子具有不大于原模的投射维数。
\end{corollary}
\begin{proof}
若 \(k\) 有长度 \(n\) 的投射消解，用它计算第二变量便得
\(\operatorname{Tor}_{n+1}^R(M,k)=0\)，由\cref{b3:thm:local-pd-vanishing}得到结论。若 \(\operatorname{pd}_Rk=\infty\)，不等式没有额外内容。对直和因子，张量复形及其同调保留有限直和；原模的该阶 Tor 为零时，各因子也为零，再应用同一判据。
\end{proof}

因此，只要剩余域的投射维数有限，它就给出本环所有有限生成模的投射维数的统一上界。

\ProseParagraphBreak
对消解逐项模掉 \(x\)，就是张量 \(R/(x)\)，而张量一般不能保留左端正合性。下面要求 \(x\) 同时在环和系数模上为非零因子，正是为了保证取商后的增广消解仍正合。

\begin{lemma}\label{b3:lem:regular-element-resolution-basechange}
设 \((R,\mathfrak m)\) 为交换 Noether 局部环，\(M\ne0\) 为有限生成的 \(R\) 模，\(x\in\mathfrak m\) 在 \(R\) 和 \(M\) 上均为非零因子。令 \(B=R/(x)\)。把 \(M\) 的极小自由消解模掉 \(x\)，得到 \(M/xM\) 的极小 \(B\) 自由消解。因此
\[
\operatorname{pd}_B(M/xM)=\operatorname{pd}_RM.
\]
\end{lemma}
\begin{proof}
因为 \(x\) 在 \(R\) 上为非零因子，\(B\) 有长度一的 \(R\) 自由消解
\[
0\longrightarrow R\xrightarrow{x}R\longrightarrow B\longrightarrow0.
\]
用这条消解及\cref{b3:lem:tor-balanced-computation}计算，得到
\[
\operatorname{Tor}_1^R(M,B)\cong\ker(x:M\rightarrow M)=0,
\qquad
\operatorname{Tor}_i^R(M,B)=0\quad(i\ge2).
\]
一次同调的消失用到了 \(x\) 在 \(M\) 上为非零因子；高次消失则来自这条消解没有更高项。因此，若 \(F\) 为 \(M\) 的极小自由消解，便有
\[
H_i(F\otimes_RB)=\operatorname{Tor}_i^R(M,B)=0\quad(i>0),
\qquad H_0(F\otimes_RB)=M\otimes_RB\cong M/xM.
\]
所以模掉 \(x\) 后的增广复形仍正合。各项是有限秩自由 \(B\) 模，微分矩阵的元素仍落在 \(B\) 的极大理想 \(\mathfrak m/(x)\) 中，故消解仍极小。Nakayama 引理给出 \(M/xM\ne0\)，而 \(F_i/xF_i\ne0\) 当且仅当 \(F_i\ne0\)。由\cref{b3:prop:pd-minimal-length}，两个极小消解的长度相同，包括无穷情形，得到所述等式。
\end{proof}

仍设 \((R,\mathfrak m)\) 为交换 Noether 局部环，\(x\in\mathfrak m\) 在 \(R\) 上为非零因子，若改取 \(M=B=R/(x)\)，则 \(xM=0\)。长度一的 \(R\) 自由消解的微分为 \(x\)，故它极小，给出 \(\operatorname{pd}_RB=1\)；而 \(B\) 在自身上自由，\(\operatorname{pd}_BB=0\)。此时 \(\operatorname{Tor}_1^R(B,B)=B\ne0\)，所以模掉 \(x\) 后的消解不再正合。共同非零因子的条件在这里不能删去。

\ProseParagraphBreak
自由项的秩、剩余域上的同调与消解的终止性，现在都可以在极小消解中读取。\cref{b3:prop:minimal-resolution-betti,b3:cor:betti-gap-termination,b3:prop:pd-minimal-length}给出的对应集中列于\cref{b3:tab:minimal-resolution-dictionary}。

\begin{center}
\begin{minipage}{\textwidth}
\makeatletter
\def\@captype{table}
\makeatother
\centering
\caption[极小自由消解的同调解释]{极小自由消解的同调解释。设 \((R,\mathfrak m)\) 为交换 Noether 局部环，\(k=R/\mathfrak m\)，\(M\ne0\) 为有限生成的 \(R\) 模，\(F_\bullet\rightarrow M\rightarrow0\) 为其极小自由消解。下表中 \(i\ge0\)，关系模指极小满射 \(F_0\rightarrow M\) 的核。}
\label{b3:tab:minimal-resolution-dictionary}
\begin{tabularx}{\textwidth}{@{}>{\raggedright\arraybackslash}p{0.36\textwidth}>{\raggedright\arraybackslash}X@{}}
\toprule
极小消解数据 & 同调解释或数值意义\\
\midrule
\(\operatorname{rank}_R F_i\) & \(\beta_i^R(M)\)\\
\(F_i/\mathfrak mF_i\) & \(\operatorname{Tor}_i^R(M,k)\)\\
\(\operatorname{Hom}_R(F_i,k)\) & \(\operatorname{Ext}_R^i(M,k)\)\\
\(F_0\) 的秩 & \(M\) 的最少生成元数\\
\(F_1\) 的秩 & 关系模的最少生成元数\\
某一阶 \(\beta_i^R(M)=0\) & 等价于 \(F_i=0\)、\(\operatorname{Tor}_i^R(M,k)=0\) 或 \(\operatorname{Ext}_R^i(M,k)=0\)，并迫使 \(F_j=0\) 对所有 \(j\ge i\) 成立\\
最后非零项为 \(F_r\) & \(\operatorname{pd}_RM=r\)\\
每个 \(F_i\) 都非零 & \(\operatorname{pd}_RM=\infty\)，没有最后非零项\\
\bottomrule
\end{tabularx}
\end{minipage}
\end{center}

\begin{proposition}\label{b3:prop:pd-test-module-counterexamples}
设 \(k\) 为域，\(R=k[x,y]_{(x,y)}\)、\(M=R/(x)\)。则
\[
\operatorname{Tor}_1^R(M,R/(y))=0,\qquad
\operatorname{Tor}_1^R(M,k)\cong k,\qquad \operatorname{pd}_RM=1.
\]
另令 \(S=k[\epsilon]/(\epsilon^2)\)，把 \(k=S/(\epsilon)\) 视为 \(S\) 模。则
\[
\begin{gathered}
\operatorname{Ext}_S^i(k,S)=0\quad(i>0),\qquad
\operatorname{Ext}_S^i(k,k)\cong k\quad(i\ge0),\\
\operatorname{pd}_Sk=\infty.
\end{gathered}
\]
因此投射维数的局部判据中，不能将剩余域任意换成别的测试模。
\end{proposition}
\begin{proof}
\(x\) 在 \(R\) 上为非零因子，故
\[
0\longrightarrow R\xrightarrow{x}R\longrightarrow M\longrightarrow0
\]
是自由消解。\(x\in(x,y)R\) 使它极小，最后非零自由项在次数一，所以 \(\operatorname{pd}_RM=1\)。张量 \(R/(y)\cong k[x]_{(x)}\) 后，乘法 \(x\) 仍单射，一次 Tor 为零；张量剩余域 \(k\) 后，微分却为零，一次 Tor 等于 \(k\)。非零的测试模 \(R/(y)\) 没有检测出这一项。

在 \(S\) 中，使用\secref{b3:sec:1701}给出的周期消解：每一项为 \(S\)，全部正次微分为乘 \(\epsilon\)。它们的核与像都为 \((\epsilon)\)，所以消解正合；各项秩为一，微分又属于极大理想，故它极小。由\cref{b3:prop:pd-minimal-length}，\(\operatorname{pd}_Sk=\infty\)。映入 \(S\) 后，上链复形的微分仍为乘 \(\epsilon\)，其正次的核与像相等，故正次 Ext 为零。映入 \(k\) 后，微分全部为零，每个次数的 Ext 都为 \(k\)。因此即使第二变量取环本身，也可能看不见无限消解中的全部正次项。
\end{proof}

\begin{corollary}\label{b3:cor:completion-betti-pd}
设 \((R,\mathfrak m)\) 为交换 Noether 局部环，\(k=R/\mathfrak m\) 为剩余域，\(M\) 为有限生成的 \(R\) 模。令 \(\widehat R\) 为 \(R\) 的 \(\mathfrak m\) 进完备化，\(\widehat M=\widehat R\otimes_RM\)。则对所有整数 \(i\ge0\)，有
\[
\beta_i^{\widehat R}(\widehat M)=\beta_i^R(M),
\qquad \operatorname{pd}_{\widehat R}\widehat M=\operatorname{pd}_RM.
\]
投射维数的等式包括 \(\infty\) 的情形；\(M=0\) 时两边均按约定为 \(-\infty\)，全部 Betti 数为零。
\end{corollary}
\begin{proof}
\(M=0\) 时直接成立，以下设 \(M\ne0\)。由\cref{b3:thm:completion-tensor}，所定义的 \(\widehat M\) 也是 \(M\) 的 \(\mathfrak m\) 进完备化。\cref{b3:thm:completion-flat}给出 \(R\to\widehat R\) 忠实平坦，所以 \(\widehat M\ne0\)；它又是有限生成的 \(\widehat R\) 模。由\cref{b3:prop:completed-local-ring}，\(\widehat R\) 是 Noether 局部环，其极大理想为 \(\mathfrak m\widehat R\)，剩余域自然等于 \(k\)。

取 \(M\) 的极小自由消解 \(F_\bullet\)。平坦性使 \(F_\bullet\otimes_R\widehat R\to\widehat M\to0\) 仍正合，各项是同秩的有限自由模。原微分的矩阵系数在 \(\mathfrak m\) 中，变基后仍在 \(\mathfrak m\widehat R\) 中，所以这也是极小自由消解。再张量共同的剩余域，得到
\[
(F_i\otimes_R\widehat R)\otimes_{\widehat R}k
\cong F_i\otimes_Rk.
\]
由\cref{b3:prop:minimal-resolution-betti}，两边的 \(k\) 维数就是各自的 Betti 数，故逐次相等。这两个极小消解的非零项出现在完全相同的次数；再用\cref{b3:prop:pd-minimal-length}，得到有限或无限的投射维数相等。
\end{proof}

% !TeX root = ../../Book3.tex
% 纲要标题：### 17.3 Auslander–Buchsbaum 公式
% 纲要位置：计划纲要.md，第 2504 行。

\section{Auslander–Buchsbaum 公式}
\label{b3:sec:1703}

在有限长度、各项均为有限秩自由模的极小消解中，末端非零合冲模是自由模，因而具有底环的深度。逐次比较消解中的短正合列，可以把这一终点与原模的深度联系起来，得到 Auslander–Buchsbaum 公式。

% 纲要内容（仅作写作计划，尚非教材正文）：
%
% * Noether 局部环及非零有限生成模的假设
% * 有限投射维数条件
% * 合并定理与证明思路；投射维数为一时用极小矩阵的 Ext 零作用，高次归纳调用深度引理
% * 公式 pd(M)+depth(M)=depth(R) 及有限投射维数下的合冲递推
%

\subsection{定理与证明思路}
\label{b3:subsec:1703:01}
\label{b3:subsec:1703:02}

深度问“能连续除去多少个非零因子”，投射维数问“自由消解需要多少步”。这两种计数一般没有相加公式；只有消解确实终止，末端自由项才能把两者联系起来。

\begin{theorem}[Auslander–Buchsbaum]\label{b3:thm:auslander-buchsbaum}
设 \((R,\mathfrak m)\) 为交换 Noether 局部环，\(k=R/\mathfrak m\) 为其剩余域，\(M\ne0\) 为有限生成的 \(R\) 模，且 \(\operatorname{pd}_RM<\infty\)。则
\begin{equation}\label{b3:eq:auslander-buchsbaum}
\operatorname{pd}_RM+\operatorname{depth}_RM=\operatorname{depth}R.
\end{equation}
\end{theorem}

这称为\term[auslander-buchsbaum-gongshi]{Auslander–Buchsbaum 公式}[Auslander--Buchsbaum formula]。局部性使极小微分的系数同时属于一个极大理想；Noether 性与有限生成使各合冲模都有有限的极小生成组，而有限投射维数保证消解终止；\(M\ne0\) 则保证深度是有限整数。证明的方法可对照\cite[Theorem 19.1, p. 155]{Matsumura1989CommutativeRingTheory}。

\ProseParagraphBreak

若 \(M\) 本身是非零有限自由模，它和 \(R\) 具有同样的正则序列，所以公式立即成立。若 \(\operatorname{pd}_RM=1\)，极小消解为
\[
0\longrightarrow F_1\xrightarrow{u}F_0\longrightarrow M\longrightarrow0.
\]
关键不只是两端自由，而是 \(u\) 的矩阵元素都在 \(\mathfrak m\) 中。它们在 \(\operatorname{Ext}_R^i(k,R)\) 上作用为零，这正使深度下降一。

\ProseParagraphBreak
这一零化作用正是\cref{b3:lem:ext-basic-properties}的第二项在第一变量为 \(k\) 时的应用：\(\mathfrak m\) 零化 \(k\)，故也零化全部 \(\operatorname{Ext}_R^i(k,R)\)。

\subsection{Auslander–Buchsbaum 公式的证明}
\label{b3:subsec:1703:03}

\begin{proof}[\cref{b3:thm:auslander-buchsbaum}的证明]
记 \(r=\operatorname{pd}_RM\)、\(t=\operatorname{depth}R\)，对 \(r\) 归纳。\(r=0\) 时，\(M\) 为非零有限生成投射模，因而平坦；由\cref{b3:lem:finite-flat-local-free}，它在局部环上自由，故与 \(R\) 具有同样深度。

设 \(r=1\)，取上一小节的极小消解，其中 \(F_1\ne0\)。应用 \(\operatorname{Hom}_R(k,-)\) 的长正合列。由\cref{b3:lem:ext-basic-properties}，映射
\[
\operatorname{Ext}_R^i(k,F_1)\longrightarrow\operatorname{Ext}_R^i(k,F_0)
\]
为零：它是 \(u\) 的同一矩阵作用于若干个 \(\operatorname{Ext}_R^i(k,R)\)。若 \(t=0\)，则 \(\operatorname{Hom}_R(k,F_1)\ne0\)，但单射 \(F_1\to F_0\) 使它必须单射地映入 \(\operatorname{Hom}_R(k,F_0)\)，与该映射为零矛盾。故 \(t\ge1\)。

对 \(i<t-1\)，长正合列两侧的自由项 Ext 均为零，故 \(\operatorname{Ext}_R^i(k,M)=0\)。在次数 \(t-1\)，得到
\[
\operatorname{Ext}_R^{t-1}(k,M)
\cong\operatorname{Ext}_R^t(k,F_1)\ne0.
\]
\cref{b3:thm:depth-ext}给出 \(\operatorname{depth}M=t-1\)。

若 \(r\ge2\)，取极小满射的核 \(0\to K\to F_0\to M\to0\)。由\cref{b3:prop:pd-minimal-length}，\(K\ne0\) 且投射维数为 \(r-1\)。归纳得
\(s\coloneqq\operatorname{depth}K=t-r+1<t=\operatorname{depth}F_0\)。若 \(s=0\)，则 \(K\to F_0\) 的单射在 \(\operatorname{Hom}_R(k,-)\) 下给出非零模到零模的单射，矛盾。因此 \(s\ge1\)。现在对该短正合列应用\cref{b3:thm:depth-lemma}的第二条附加等式，得到
\[
\operatorname{depth}M=\operatorname{depth}K-1=s-1=t-r.
\]
归纳完成。
\end{proof}

把证明中的递推放回整个极小消解，可以看出每一步如何联系两种计数。设 \((R,\mathfrak m)\) 为交换 Noether 局部环，\(M\ne0\) 为有限生成的 \(R\) 模，\(r=\operatorname{pd}_RM<\infty\)、\(t=\operatorname{depth}R\)，并取 \(M\) 的极小自由消解。在其中令 \(\Omega^0M=M\)、\(\Omega^{j+1}M=\ker(F_j\to\Omega^jM)\)，则
\[
\begin{gathered}
0\longrightarrow\Omega^{j+1}M\longrightarrow F_j
\longrightarrow\Omega^jM\longrightarrow0
\qquad(0\le j<r),\\[2pt]
\operatorname{pd}_R\Omega^{j+1}M
=\operatorname{pd}_R\Omega^jM-1,
\qquad
\operatorname{depth}_R\Omega^{j+1}M
=\operatorname{depth}_R\Omega^jM+1,\\[2pt]
\Omega^rM\cong F_r\ne0,
\qquad \operatorname{depth}_R\Omega^rM=t.
\end{gathered}
\]
第一行中的范围也适用于第二行：在到达非零有限自由模之前，投射维数每次减少一，深度每次增加一；末端的深度等于 \(t\)，因而起点深度为 \(t-r\)。这也给出非零有限生成模在投射维数有限时的上界 \(\operatorname{pd}_RM\le\operatorname{depth}R\)。

% !TeX root = ../../Book3.tex
% 纲要标题：### 17.4 投射维数的计算与有限性条件
% 纲要位置：计划纲要.md，第 2511 行。

\section{投射维数的计算与有限性条件}
\label{b3:sec:1704}

正则序列给出可直接写下的有限消解，而有些小商环的关系会一直重复，产生无限消解。下面通过两类算例同时核对投射维数、Betti 数与深度，避免只凭方程个数估计消解长度。

% 纲要内容（仅作写作计划，尚非教材正文）：
%
% * 商环、理想及合冲模的计算
% * 不能对无限投射维数直接套公式
% * 为正则局部环的同调刻画准备工具
%
% ---
%

\subsection{商环、理想与合冲模的计算}
\label{b3:subsec:1704:01}

若 \(x_1,\ldots,x_c\in\mathfrak m\) 是交换 Noether 局部环 \(R\) 上的正则序列，则\cref{b3:thm:koszul-acyclic}给出商模的长度为 \(c\)、各项有限秩自由的 Koszul 消解。因为各 \(x_i\in\mathfrak m\)，它极小，末项为 \(R\)，所以
\[
\operatorname{pd}_R R/(x_1,\ldots,x_c)=c,\qquad
\beta_i^R\bigl(R/(x_1,\ldots,x_c)\bigr)=\binom ci.
\]
Auslander–Buchsbaum 公式给出商模深度下降 \(c\)，与逐次除去非零因子的深度计算一致。

\ProseParagraphBreak
设 \(k\) 为域，取 \(R=k[x,y]_{(x,y)}\)、\(I=(x,y)\)。上述 Koszul 消解去掉最后的 \(R\to k\)，得到
\[
0\longrightarrow R\xrightarrow{\binom{-y}{x}}R^2\longrightarrow I\longrightarrow0.
\]
因此 \(\operatorname{pd}_RI=1\)、\(\operatorname{depth}_RI=1\)。理想 \(I\) 是整环中的无挠模，却不是自由模；无挠只能排除单个非零元素的零化，不能保证深度与环一样大。

\begin{corollary}\label{b3:cor:syzygy-depth-finite-pd}
设 \((R,\mathfrak m)\) 为交换 Noether 局部环，\(M\ne0\) 为有限生成的 \(R\) 模且 \(r=\operatorname{pd}_RM<\infty\)。令 \(\Omega^0M=M\)，\(\Omega^jM\) 为 \(M\) 的极小自由消解中的第 \(j\) 个合冲模。对每个整数 \(0\le j\le r\)，\(\Omega^jM\ne0\)，并满足
\[
\operatorname{pd}_R\Omega^jM=r-j,\qquad
\operatorname{depth}_R\Omega^jM=\operatorname{depth}_RM+j.
\]
\end{corollary}
\begin{proof}
截去极小消解的前 \(j\) 项得到 \(\Omega^jM\) 的极小消解，长度为 \(r-j\)，且末项仍非零。分别对 \(M\) 和 \(\Omega^jM\) 应用\cref{b3:thm:auslander-buchsbaum}，相减得到深度等式。
\end{proof}

\subsection{有限投射维数假设的必要性}
\label{b3:subsec:1704:02}

设 \(k\) 为域。在 \(R=k[\epsilon]/(\epsilon^2)\) 上，\(\operatorname{depth}R=0\)，剩余域 \(k\) 的深度也为零。然而\secref{b3:sec:1701}已算出 \(k\) 的每个 Betti 数都为一，故 \(\operatorname{pd}_Rk=\infty\)。不能把两个深度之差为零解释为 \(k\) 投射。

\ProseParagraphBreak
更有几何意味的例子是 \(R=k[x,y]_{(x,y)}/(xy)\) 与 \(M=R/(x)\)。\(R\) 中有
\(\operatorname{Ann}(x)=(y)\)、\(\operatorname{Ann}(y)=(x)\)：在多项式整环中约去 \(x\) 或 \(y\) 即得，局部化保持这些核。因此
\[
\cdots\xrightarrow{x}R\xrightarrow{y}R\xrightarrow{x}R
\longrightarrow M\longrightarrow0
\]
为极小自由消解，\(M\) 的投射维数无限。另一方面 \(M\cong k[y]_{(y)}\)，故 \(\operatorname{depth}_RM=1\)。\(R\) 的维数为一，而 \(x+y\) 不属于两个关联素理想 \((x),(y)\)，是非零因子，故 \(\operatorname{depth}R=1\)。深度相等仍不表示自由。

\subsection{正则局部环的同调刻画初步}
\label{b3:subsec:1704:03}

\begin{corollary}\label{b3:cor:finite-pd-maximal-depth-free}
设 \((R,\mathfrak m)\) 为交换 Noether 局部环，\(M\ne0\) 为有限生成的 \(R\) 模。若 \(M\) 投射维数有限且 \(\operatorname{depth}_RM=\operatorname{depth}R\)，则 \(M\) 自由。特别地，深度为零的这种局部环上，每个具有有限投射维数的有限生成模都自由。
\end{corollary}
\begin{proof}
\cref{b3:thm:auslander-buchsbaum}使投射维数为零，因而模投射；局部有限投射模由\cref{b3:thm:finite-presentation-flat-projective}和\cref{b3:lem:finite-flat-local-free}自由。深度为零时，公式右边为零而左边两个整数均非负，故仍然如此。
\end{proof}

对交换 Noether 局部环 \((R,\mathfrak m)\)，剩余域 \(k=R/\mathfrak m\) 尤其敏感。若 \(\operatorname{pd}_Rk\) 有限，本章已经说明所有有限生成的 \(R\) 模的投射维数都有统一上界，且 \(\operatorname{pd}_Rk=\operatorname{depth}R\)。下一章将结合极大理想的生成元，证明这个同调条件还迫使 \(\dim R=\operatorname{edim}R\)。


\begin{chaptersummary}
极小消解把局部环上有限生成模的生成关系变成一列由模自身决定的自由秩。剩余域使微分消失，所以任一正次 Betti 数为零就表示消解已经终止。对一般消解，这样直接读取自由秩的方法不成立。

\ProseParagraphBreak
有限投射维数使自由消解具有末端，Auslander–Buchsbaum 公式由此把消解长度和深度相加为环的深度。公式在投射维数为一时依赖极小矩阵对剩余域 Ext 的零作用，较长情形再逐步取合冲模。周期消解说明深度相等不能单独推出自由；下一章把剩余域的消解终止与环的正则性联系起来。
\end{chaptersummary}

\BookExercises

\begin{exercise}\label{b3:ex:17-contractible-summand}
设 \((R,\mathfrak m)\) 为交换 Noether 局部整环，\(0\ne x\in\mathfrak m\)。令 \(\epsilon:R^2\to R/(x)\) 为 \(\epsilon(a,b)=a\bmod(x)\)。验证
\[
0\longrightarrow R^2\xrightarrow{\operatorname{diag}(x,1)}R^2
\xrightarrow{\epsilon}R/(x)\longrightarrow0
\]
是长度一的自由消解，并从中分离出极小消解。说明为什么不能把该矩阵的阶数作为 \(\beta_1^R(R/(x))\)。
\end{exercise}
\begin{solution}
矩阵单射，因为 \(x\) 为非零因子；其余核由第一坐标模 \(x\) 给出，故余核为 \(R/(x)\)。此消解是 \(0\to R\xrightarrow{x}R\to R/(x)\to0\) 与零增广的 \(R\xrightarrow{1}R\) 的直和。后一复形可缩且没有同调，却把两项自由秩各增加一。极小消解的两个秩都是一，故 \(\beta_0=\beta_1=1\)。
\end{solution}

\begin{exercise}\label{b3:ex:17-dvr-betti}
设 \(R\) 为离散赋值环，统一化元为 \(\pi\)，\(r,s\ge0\) 为整数，\(a,b,a_1,\ldots,a_s\) 为正整数。令
\[
M=R^r\oplus\bigoplus_{j=1}^sR/(\pi^{a_j}).
\]
求 \(M\) 的 Betti 数，再求
\[
\operatorname{Tor}_1^R(R/(\pi^a),R/(\pi^b)),\qquad
\operatorname{Ext}_R^1(R/(\pi^a),R/(\pi^b)).
\]
\end{exercise}
\begin{solution}
把各循环模的 \(0\to R\xrightarrow{\pi^{a_j}}R\to R/(\pi^{a_j})\to0\) 直和，再接上自由项，得到极小消解。因此 \(\beta_0=r+s\)、\(\beta_1=s\)，其余为零。
记 \(N=R/(\pi^b)\)。Tor 是 \(N\xrightarrow{\pi^a}N\) 的核，即 \(\pi^{\max(b-a,0)}R/\pi^bR\)，同构于 \(R/(\pi^{\min(a,b)})\)。Ext 是同一乘法的余核 \(N/\pi^aN\)，也同构于 \(R/(\pi^{\min(a,b)})\)。两个同构的取得方式不同：一个由子模，一个由商模得到。
\end{solution}

\begin{exercise}\label{b3:ex:17-embedded-point-resolution}
设 \(k\) 为域，\(R=k[x,y]_{(x,y)}\)、\(M=R/(x^2,xy)\)。构造其极小自由消解，计算投射维数与深度，并求其基座 \(\operatorname{Soc}_RM\)。
\end{exercise}
\begin{solution}
消解为
\[
0\to R\xrightarrow{\binom{-y}{x}}R^2
\xrightarrow{(x^2\quad xy)}R\to M\to0.
\]
若 \(ax^2+bxy=0\)，约去 \(x\) 后有 \(ax+by=0\)。模 \(x\) 得 \(b\in(x)\)，写 \(b=cx\)，则 \(a=-cy\)，故核恰为所示像；左箭头单射。微分元素均在极大理想中，消解极小，所以投射维数二。\(R\) 的深度二，公式给 \(\operatorname{depth}M=0\)。非零元 \(\bar x\) 被 \(x,y\) 零化。反之，若 \(\mathfrak m\bar f=0\)，把 \(\bar f\) 映到 \(M/(\bar x)\cong k[y]_{(y)}\)，乘 \(y\) 的单射性使其像为零，所以 \(\bar f\in R\bar x\)。又 \(\mathfrak m\bar x=0\)，故 \(R\bar x=k\bar x\)，从而 \(\operatorname{Soc}_RM=k\bar x\)。
\end{solution}

\begin{exercise}\label{b3:ex:17-square-maximal-ideal}
设 \(k\) 为域，在 \(R=k[x,y]_{(x,y)}\) 中，求 \(I=(x,y)^2\) 以及 \(R/I\) 的极小自由消解和 Betti 数。
\end{exercise}
\begin{solution}
\(I\) 的生成元为 \(x^2,xy,y^2\)，关系由 \((-y,x,0)\)、\((0,-y,x)\) 生成。确实，若 \(ax^2+bxy+cy^2=0\)，模 \(x\) 得 \(c=xv\)，约去 \(x\) 得 \(ax+(b+vy)y=0\)，再得 \(b+vy=xu\)、\(a=-uy\)。因此关系向量为 \(u(-y,x,0)+v(0,-y,x)\)。两向量线性独立，因为第一坐标先使 \(u=0\)，第三坐标再使 \(v=0\)。所以
\[
0\to R^2\xrightarrow{\left(\begin{smallmatrix}-y&0\\x&-y\\0&x\end{smallmatrix}\right)}R^3\to I\to0
\]
极小，\(I\) 的 Betti 数为 \((3,2)\)。接上 \(I\to R\to R/I\) 得商的极小消解，Betti 数为 \((1,3,2)\)。
\end{solution}

\begin{exercise}\label{b3:ex:17-tor-choice-matters}
设 \(k\) 为域，令 \(R=k[x,y]_{(x,y)}\)、\(M=R/(x)\)，其中 \(k\) 按剩余域视为 \(R\) 模。分别计算
\(\operatorname{Tor}_1^R(M,R/(y))\)、\(\operatorname{Tor}_1^R(M,M)\) 和 \(\operatorname{Tor}_1^R(M,k)\)。
\end{exercise}
\begin{solution}
以 \(0\to R\xrightarrow{x}R\to M\to0\) 计算。张量 \(R/(y)\) 后乘 \(x\) 单射，所以第一个 Tor 为零。张量 \(M\) 后乘 \(x\) 为零，第二个 Tor 为 \(M\)。张量 \(k\) 后也为零，第三个 Tor 为 \(k\)。尽管第一个测试结果为零，\(M\) 的投射维数仍为一。
\end{solution}

\begin{exercise}\label{b3:ex:17-dual-number-ext}
设 \(k\) 为域，令 \(R=k[\epsilon]/(\epsilon^2)\)，其剩余域自然识别为 \(k\)。求全部 \(\operatorname{Ext}_R^i(k,k)\) 与 \(\operatorname{Ext}_R^i(k,R)\)，并说明只以 \(R\) 为第二变量测试可能漏掉什么。
\end{exercise}
\begin{solution}
用正文的周期极小消解。映入 \(k\) 后微分全零，故每个 \(i\ge0\) 都有 \(\operatorname{Ext}_R^i(k,k)\cong k\)。映入 \(R\) 后仍为连续乘 \(\epsilon\) 的复形；正次数的核、像都为 \((\epsilon)\)，所以这些 Ext 全零，而零次 Hom 为 \((\epsilon)\cong k\)。后一个消失没有使 \(k\) 投射维数有限，因而不能将局部判据中的测试剩余域擅自换成 \(R\)。
\end{solution}

\begin{exercise}\label{b3:ex:17-direct-sum-betti}
设 \((R,\mathfrak m)\) 为交换 Noether 局部环，\(M,N\) 为有限生成的 \(R\) 模。证明对每个整数 \(i\ge0\)，有
\(\beta_i^R(M\oplus N)=\beta_i^R(M)+\beta_i^R(N)\)。另设 \(k\) 为域，求
\(k\oplus R/(x)\) 在 \(R=k[x,y]_{(x,y)}\) 上的 Betti 数与深度，其中 \(k\) 为 \(R\) 的剩余域模。
\end{exercise}
\begin{solution}
两极小消解的有限直和仍正合、自由且极小，所以各自由秩相加。例子中两列 Betti 数为 \((1,2,1)\) 与 \((1,1)\)，和为 \((2,3,1)\)。因此投射维数二，深度为零；也可由直和含一个被极大理想零化的非零元素直接得到深度零。
\end{solution}

\begin{exercise}\label{b3:ex:17-rank-euler}
设 \(R\) 为交换 Noether 局部整环，\(M\) 为投射维数有限的有限生成模。证明
\(\operatorname{rank}_RM=\sum_i(-1)^i\beta_i^R(M)\)。用它核对上题和平方极大理想的计算。
\end{exercise}
\begin{solution}
把有限长度的极小自由消解张量到分式域 \(K\)。局部化正合，得到有限维向量空间的正合列；将每个核的维数在相邻两项中抵消，得到末端维数等于自由秩的交错和。上题两个商模都是挠模，和的秩为零，\(2-3+1=0\)。理想 \((x,y)^2\) 的秩为一，\(3-2=1\)。
\end{solution}

\begin{exercise}\label{b3:ex:17-completion-betti}
设 \((R,\mathfrak m)\) 为交换 Noether 局部环，\(k=R/\mathfrak m\) 为其剩余域，\(M\) 为有限生成的 \(R\) 模。证明完备化保持全部 Betti 数和投射维数，包括投射维数无穷的情形。
\end{exercise}
\begin{solution}
\(\widehat R\) 忠实平坦，极大理想为 \(\mathfrak m\widehat R\)。把极小自由消解张量 \(\widehat R\) 后仍正合，微分仍落在新极大理想中，且自由项的秩不变；其末端为 \(M\otimes_R\widehat R\cong\widehat M\)。因此它是完备化模的极小消解。两列 Betti 数相同，其非零次数的上确界也相同。
\end{solution}

\begin{exercise}\label{b3:ex:17-localization-pd-drop}
设 \(k\) 为域，\(R=k[x,y]_{(x,y)}\)、\(I=(x,y)\)。比较 \(I\) 在 \(R\) 上及 \(I_{(x)}\) 在 \(R_{(x)}\) 上的投射维数，并解释为什么局部化极小消解后不一定仍极小。
\end{exercise}
\begin{solution}
正文给出 \(\operatorname{pd}_RI=1\)。在 \((x)\) 处 \(y\) 成为单位，故 \(I_{(x)}=R_{(x)}\)，投射维数零。原关系矩阵 \(\binom{-y}{x}\) 局部化后出现单位项，可以拆掉一个可缩的二项自由复形；微分不再全落在 \((x)R_{(x)}\) 中。
\end{solution}

\begin{exercise}\label{b3:ex:17-basechange-hypothesis}
设 \(k\) 为域，在 \(R=k[x,y]_{(x,y)}\) 中取 \(B=R/(x)\)。分别对 \(M=R/(y)\) 与 \(N=R/(x)\) 比较原环和商环上的消解，说明共同非零因子条件的作用。
\end{exercise}
\begin{solution}
\(x\) 在 \(R\) 和 \(M\) 上均单射，故 \(M/xM=k\) 的 \(B\) 消解为 \(0\to B\xrightarrow{y}B\to k\to0\)，投射维数仍为一。
但 \(x\) 在 \(N\) 上为零。\(N\) 原来的消解模 \(x\) 后得到零微分 \(B\xrightarrow{0}B\)，一次同调为 \(B\)，不是消解。\(N/xN=B\) 在 \(B\) 上自由，投射维数从一变成零。
\end{solution}

\begin{exercise}\label{b3:ex:17-node-depth-one}
设 \(k\) 为域，\(R=k[x,y]_{(x,y)}/(xy)\)。对 \(M=R/(x)\)，确定各个极小合冲模的同构类型，核对它们的深度。
\end{exercise}
\begin{solution}
第一合冲模为 \((x)\cong R/(y)\)，第二合冲模为 \((y)\cong R/(x)\)，此后交替。前一同构来自乘 \(x\) 的核为 \((y)\)，后一同理。这两个商分别是 \(k[x]_{(x)}\) 和 \(k[y]_{(y)}\)，作为 \(R\) 模深度都为一。无限消解中的深度达到环的深度后保持不变，并不迫使消解终止。
\end{solution}

\begin{exercise}\label{b3:ex:17-nonfinite-warning}
设 \(p\) 为素数，令 \(R=\mathbb Z_{(p)}\)、\(K=\mathbb Q\)。证明 \(K\) 平坦但不是有限生成 \(R\) 模，且 \(K/pK=0\)。说明为什么不能用“剩余域张量为零”在此推出 \(K=0\)。
\end{exercise}
\begin{solution}
\(K\) 是 \(R\) 的局部化，故平坦。若由有限个分式生成，取它们分母中 \(p\) 的最高次数 \(a\)，生成模将包含于 \(p^{-a}R\)，但 \(p^{-a-1}\in K\) 不在其中，矛盾。由于 \(p\) 在 \(K\) 中可逆，\(pK=K\)。Nakayama 引理要求有限生成；这里既没有有限极小生成组，也不能用 \(\beta_0=0\) 的有限生成模结论。
\end{solution}

\begin{exercise}\label{b3:ex:17-unit-cancellation}
设 \((R,\mathfrak m)\) 为交换 Noether 局部环，\(M\) 为有限生成的 \(R\) 模，给定 \(M\) 的一个有限长度、各项有限秩自由的消解 \(F_\bullet\)。令 \(b_i=\operatorname{rank}_RF_i\)。证明：从中逐次消去单位项得到极小消解时，拆出的二项复形 \(R\xrightarrow{1}R\) 的总数为
\[
\dfrac{1}{2}\sum_{i\ge0}\bigl(b_i-\beta_i^R(M)\bigr),
\]
因而不依赖消去次序。
\end{exercise}
\begin{solution}
由\cref{b3:cor:finite-resolution-splitting}，有限步消去后，原消解分解为极小消解与若干二项复形的直和。极小消解在次数 \(i\) 的自由秩为 \(\beta_i^R(M)\)，每个拆出的二项复形对总自由秩恰贡献二。因此若拆出 \(c\) 个二项复形，则
\[
\sum_{i\ge0}b_i=\sum_{i\ge0}\beta_i^R(M)+2c.
\]
两端的和都只有有限多个非零项，解出 \(c\) 即得题中公式，也说明消去次序不影响总数。
\end{solution}
